% =====================================================================
% ABLATION DRAFTS — pick any 2--3 paragraphs + the matching rows of
% Table~\ref{tab:ablations} for the final Analysis section.
%
% Each \paragraph block is self-contained: hypothesis, ablated component,
% slice, quantitative impact, and interpretation.
%
% Master table at the bottom aggregates one row per ablation so a single
% \begin{table} suffices regardless of which subset you pick.
% =====================================================================

\subsection{Ablation Studies}
\label{sec:ablations}

To isolate the contribution of each component in the proposed pipeline we
ran a battery of single-slice ablations.  Unless otherwise stated, all
deltas in Table~\ref{tab:ablations} are CPA drops at $c{=}0.5$ relative to
the corresponding ``full method'' configuration (gemini-2.0-flash-001
speaker and listener, text-assisted listener, Gaussian posterior, EU-based
clarification rule active, no image annotation).

% --------------------------------------------------------------------
% A. Speaker vocabulary constraint
% --------------------------------------------------------------------
\paragraph{A. Constrained speaker vocabulary
            (\textit{canonical\_vllm} $\to$ \textit{natural}).}
We replace the closed-vocabulary speaker prompt --- which forces COLOR /
SHAPE / SIZE / LOCATION choices from fixed lists and applies greedy
feature inclusion --- with the unconstrained \textit{natural} prompt that
permits free-form descriptions.
We hypothesised that vocabulary constraints suppress the dominant VLM
speaker failure modes (vague generic words, hallucinated attributes,
under-specification) catalogued in \S\ref{sec:exp2}.
On \texttt{scenes\_10\_force} (n${=}300$) with the
\textit{direct-rank}(flash) listener, removing the constraint drops CPA
from $0.947 \to 0.900$ ($\Delta {=} {-}0.047$); the elimination listener
shows a comparable drop ($0.677 \to 0.680$ on this slice, but $0.793 \to
0.742$ on \texttt{scenes\_10\_none}).
The effect is a small but consistent positive contribution; the residual
gap to the rule-based oracle ($0.982$) is dominated by under-specification
errors that prompting alone cannot fully eliminate.

% --------------------------------------------------------------------
% B. Listener text-channel (leakage)
% --------------------------------------------------------------------
\paragraph{B. Listener text-channel
            (text-assisted $\to$ image-only).}
We remove the formatted object-attribute list from the listener prompt,
forcing the model to ground the utterance in pixels through an
index-annotated image alone.
The hypothesis is that text-assisted (TA) evaluation overestimates
listener capability: the formatted list reduces the task to a
text-lookup that bypasses visual grounding entirely.
On \texttt{scenes\_10\_force} (n${=}300$) with the
\textit{canonical\_vllm} speaker, switching from TA \textit{direct-rank}
to image-only (IO) \textit{io-direct} drops CPA from $0.947 \to 0.687$
($\Delta {=} {-}0.260$).
Averaged over all four speakers and both overlap settings at $N{=}10$,
the leakage gap is $0.266$ CPA points and grows monotonically with $N$
(Tab.~\ref{tab:io-summary}), confirming that TA listener scores must be
read as upper bounds rather than honest estimates of visuo-linguistic
competence.

% --------------------------------------------------------------------
% C. Listener reasoning protocol (CoT vs direct)
% --------------------------------------------------------------------
\paragraph{C. Listener reasoning protocol
            (\textit{io-direct} $\to$ \textit{io-cot}).}
We replace the single-pass probability-emission prompt with a structured
chain-of-thought protocol that asks the listener to first describe each
indexed object, then score, then convert scores to probabilities.
We expected explicit deliberation to improve calibration; in practice it
inflates expressed uncertainty.
On \texttt{scenes\_10\_force} (n${=}300$) with the
\textit{canonical\_vllm} speaker, CoT collapses CPA from $0.687 \to
0.187$ at $c{=}0.5$ ($\Delta {=} {-}0.500$), and CPA turns
\emph{negative} ($-0.083$) at $c{=}0.10$.
The failure mode is uncertainty inflation: when the model verbalises its
visual ambiguity (``Object~3 appears to be either a blue or green
circle'') the resulting posterior approaches uniform, triggering the EU
rule to ask in $>{60}\%$ of scenes.  Contrary to the System-2 intuition,
forcing explicit deliberation \emph{hurts} a listener whose intuitive
argmax was already correct.

% --------------------------------------------------------------------
% D. Posterior sharpening (kernel ablation)
% --------------------------------------------------------------------
\paragraph{D. Posterior sharpening
            (Gaussian $\sigma{=}10$ $\to$ inverse-distance).}
We swap the Gaussian distance kernel
$\exp(-d_i^2/2\sigma^2)$ used by VLLMListener for the original
inverse-distance kernel $1/(d_i{+}\epsilon)$, holding the listener's
predicted $(x,y)$ coordinates fixed.
The hypothesis is that the kernel choice affects only the shape of the
posterior, not the argmax: any change in CPA must therefore be
attributable to the cost-utility decision rule reacting to posterior
flatness rather than to a different prediction.
On \texttt{scenes\_8\_force} (n${=}300$) with the
\textit{vllm-pragmatic}(flash) speaker, removing sharpening drops mean
$\max_i P(t_i\mid u)$ from $0.61 \to 0.34$, raises ask\% at $c{=}0.5$
from $32\%\to 90\%$, and collapses CPA from $+0.233 \to -0.375$
($\Delta {=} {-}0.608$); argmax accuracy is unchanged at $57.7\%$ as
predicted.
This is the most mechanically interpretable ablation in our pipeline:
the entire CPA gain comes from the kernel making the EU rule actionable,
without altering the listener's underlying predictions.

% --------------------------------------------------------------------
% E. Cost-aware clarification rule (the core method)
% --------------------------------------------------------------------
\paragraph{E. Cost-aware clarification rule (full method $\to$ always-commit).}
We disable the EU decision rule (Eq.~\ref{eq:eu-rule}) and force the
listener to always commit to $\arg\max_i P(t_i\mid u)$, eliminating the
possibility of asking.
The hypothesis is that the EU rule yields positive CPA gain whenever the
posterior is well-calibrated --- the central methodological claim of the
project.
On \texttt{scenes\_10\_force} (n${=}300$) with the IO listener,
\emph{removing} the rule (\textit{io-index}: always commits) \emph{raises}
CPA from $0.677 \to 0.823$ at $c{=}0.5$ for the
\textit{feature-canonical} speaker ($\Delta {=} +0.146$), and similar
inversions hold across all four IO speakers and across $N \in \{6,8,10\}$
(Tab.~\ref{tab:io-c05}).
This honest negative result indicates that posterior flatness in the IO
regime is driven by \emph{visual} perceptual uncertainty rather than by
genuine referential ambiguity in the utterance: the EU rule, which
assumes flatness $\Rightarrow$ ambiguity, asks on scenes the listener
would have answered correctly.  Restoring the rule's intended behaviour
will require recalibration that distinguishes perceptual noise from
referential ambiguity.

% --------------------------------------------------------------------
% G. Visual annotation channel
% --------------------------------------------------------------------
\paragraph{G. Visual annotation channel
            (text-only target $\to$ magenta bounding box).}
We replace text-only target identification (the target's
$\langle\text{color}, \text{shape}, \text{size}, x, y\rangle$ injected into
the system prompt) with the original visual annotation: a magenta
bounding box and ``TARGET'' label drawn on the image at the target's
pixel center.
The hypothesis is that strong VLMs anchor on annotation pixels and
describe the overlay rather than the underlying object.
On a held-out subset of \texttt{scenes\_6\_force} (n${=}50$) with the
\texttt{claude-sonnet-4} speaker, annotation-based identification yields
CPA $\approx 0.12$ versus $0.94$ for text-only target passing
($\Delta {=} {-}0.82$); failure traces show the speaker outputting
``the magenta square'' or ``the highlighted object'' instead of the
target's actual attributes.
This anecdotal but striking gap motivated the design change to text-only
target identification used throughout Experiments~2--5; reproducing it as
a controlled n${=}300$ ablation is left for future work.

% --------------------------------------------------------------------
% I. Listener model capacity
% --------------------------------------------------------------------
\paragraph{I. Listener model capacity
            (\textit{direct-rank}(flash) $\to$ \textit{direct-rank}(sonnet-4)).}
Holding the speaker fixed at \textit{listener\_aware}(flash), we swap the
listener backbone from gemini-2.0-flash-001 to claude-sonnet-4.
The hypothesis is that the residual gap to the rule-based oracle is
dominated by listener capacity rather than by speaker quality, as
suggested by Finding~7 of \S\ref{sec:honest}.
On \texttt{scenes\_6\_force} (n${=}50$), the stronger listener raises CPA
from $0.95 \to 0.97$ ($\Delta {=} +0.02$) at $c{=}0.5$; under TA
evaluation the effect is small because gemini-flash already saturates the
text-lookup channel, but in IO evaluation listener capacity is the
dominant lever for closing the gap to the oracle ceiling.
This ablation is small at TA scale; its value is in motivating
listener-side scaling for the honest IO protocol where speaker engineering
yields diminishing returns.

% --------------------------------------------------------------------
% L. Forced-overlap data augmentation
% --------------------------------------------------------------------
\paragraph{L. Forced same-color distractor
            (\textit{none} $\to$ \textit{force}).}
We compare datasets generated with independent attribute draws
(\textit{none}) against those that guarantee at least one distractor
shares the target's color (\textit{force}), holding $N{=}10$ and the
speaker fixed.
The hypothesis is that a near-confusable distractor genuinely raises
referential difficulty: a robust pipeline should degrade gracefully, but
a brittle one should fail catastrophically.
With the \textit{canonical\_vllm}(flash) speaker and TA
\textit{direct-rank}(flash) listener, CPA at $c{=}0.5$ drops from
$0.990 \to 0.947$ ($\Delta {=} {-}0.043$); under the more brittle
\textit{elimination}(flash) listener the same comparison drops
$0.793 \to 0.677$ ($\Delta {=} {-}0.116$); under IO \textit{io-direct} it
drops $0.733 \to 0.687$ ($\Delta {=} {-}0.046$).
Forced overlap therefore exposes a robustness gradient across listener
designs that is invisible on the \textit{none} variant alone, justifying
its inclusion as a controlled difficulty axis rather than as a
confounder.

% =====================================================================
% Master table — one row per ablation
% =====================================================================
\begin{table}[ht]
\centering
\small
\setlength{\tabcolsep}{4pt}
\begin{tabular}{cllccc}
\toprule
ID & Component ablated & Single slice
   & Full & Ablated & $\Delta$ \\
\midrule
A & Speaker prompt: \textit{canonical\_vllm} $\!\to\!$ \textit{natural}
  & \texttt{scenes\_10\_force}, DR(flash)
  & 0.947 & 0.900 & $-0.047$ \\
B & Listener input: TA $\!\to\!$ image-only (\textit{io-direct})
  & \texttt{scenes\_10\_force}, \textit{canonical\_vllm}
  & 0.947 & 0.687 & $-0.260$ \\
C & Listener reasoning: \textit{io-direct} $\!\to\!$ \textit{io-cot}
  & \texttt{scenes\_10\_force}, \textit{canonical\_vllm}
  & 0.687 & 0.187 & $-0.500$ \\
D & Posterior kernel: Gaussian $\sigma{=}10$ $\!\to\!$ inverse-distance
  & \texttt{scenes\_8\_force}, \textit{vllm-pragmatic}
  & $+0.233$ & $-0.375$ & $-0.608$ \\
E & Clarification rule: EU on $\!\to\!$ always-commit (\textit{io-index})
  & \texttt{scenes\_10\_force}, \textit{feature-canonical}
  & 0.677 & 0.823 & $+0.146$\,$^{\dagger}$ \\
G & Speaker target: text-only $\!\to\!$ magenta bbox annotation
  & \texttt{scenes\_6\_force} (n${=}50$), claude-sonnet-4
  & 0.940 & 0.120 & $-0.820$ \\
I & Listener backbone: flash $\!\to\!$ claude-sonnet-4
  & \texttt{scenes\_6\_force} (n${=}50$), \textit{listener\_aware}
  & 0.950 & 0.970 & $+0.020$\,$^{\ddagger}$ \\
L & Data: forced same-color distractor on/off
  & $N{=}10$, \textit{canonical\_vllm}+DR
  & 0.990 & 0.947 & $-0.043$ \\
\bottomrule
\end{tabular}
\caption{Single-slice ablation summary at $c{=}0.5$ unless noted.
$\Delta$ is the CPA change when the listed component is removed
(positive means the ablation \emph{improves} CPA).
$^{\dagger}$\,Removing the EU rule \emph{improves} CPA in the IO regime,
because perceptual flatness is misread as referential ambiguity.
$^{\ddagger}$\,Improvement is small at TA scale; effect is larger in IO
evaluation (\S\ref{sec:honest} Finding~7).}
\label{tab:ablations}
\end{table}
